\documentclass[12 pt, letterpaper]{hmcpset}
\usepackage{hw-preamble}

\name{Jayden Thadani}
\class{MATH 195 --- Transcendental Number Theory}
\assignment{Midterm}
\duedate{12th March 2026}

\begin{document}

\begin{problem}[Problem 1.]
	Let $\alpha, \beta \in \mathbb{C}$ and $\alpha$ transcendental. Prove that at least one of the numbers $\alpha + \beta$ and $\alpha \beta$ must be transcendental.
\end{problem}

\begin{solution}
	Note that it suffices to show this when $\beta$ is transcendental. Else, if $\beta \in \overline{\mathbb{Q}}$, since the algebraic numbers $\overline{\mathbb{Q}}$ are closed under addition, $\alpha + \beta \in \overline{\mathbb{Q}}$ implies $\alpha = (\alpha + \beta) - \beta \in \overline{\mathbb{Q}}$. Thus, if $\alpha \in \mathbb{C} \setminus \overline{\mathbb{Q}}$ and $\beta \in \overline{\mathbb{Q}}$, then $\alpha + \beta \in \mathbb{C} \setminus \overline{\mathbb{Q}}$.

	Suppose $\beta$ is transcendental. If $\alpha + \beta \in \overline{\mathbb{Q}}$ we are done. Else, write $\alpha + \beta = \gamma \in \overline{\mathbb{Q}}$. This gives $\alpha \beta = \alpha \gamma - \alpha^{2}$. If $\alpha \gamma - \alpha^{2} = \delta \in \overline{\mathbb{Q}}$ were algebraic, then we have $\alpha$ as a root of the polynomial $x^{2} - \gamma x + \delta \in \mathbb{Q}(\gamma, \delta)[x]$. We claim this implies that $\alpha$ is algebraic over $\mathbb{Q}$. This contradicts our hypothesis that $\alpha$ is transcendental, and thus, our assumption that both $\alpha + \beta$ and $\alpha \beta$ are algebraic must be false.

	First, $f(\alpha) = 0$ directly implies that $\mathbb{Q}(\alpha, \gamma, \delta) / \mathbb{Q}(\gamma, \delta)$ is a finite degree extension. Since $\mathbb{Q}(\gamma, \delta) / \mathbb{Q}$ is a finite degree extension (by choosing finite degree minimal polynomials of $\gamma, \delta$ over $\mathbb{Q}$) it follows that $\mathbb{Q}(\alpha, \gamma, \delta) / \mathbb{Q}$ is a finite degree extension, and finally, that $\mathbb{Q}(\alpha) / \mathbb{Q}$ is a finite degree extension. Thus, $\alpha$ is algebraic over $\mathbb{Q}$ (all transcendental extensions are infinite degree).
\end{solution}

\pagebreak

\begin{problem}[Problem 2.]
	Let $\mathbb{K} = \mathbb{Q}(\sqrt{2}, \sqrt{3}, \sqrt{5})$. Find a primitive element for $\mathbb{K}$. Prove your answer.
\end{problem}

\begin{solution}
	We claim $\mathbb{K} = \mathbb{Q}(\sqrt{2} + \sqrt{3} + \sqrt{5})$.

	Suppose $\alpha_{1}, \alpha_{2}$ are degree $d$ and $e$ algebraic numbers. Let $\alpha_{1, 1}, \dots, \alpha_{1, d}$ and $\alpha_{2, 1}, \dots, \alpha_{2, e}$ be all of their conjugates respectively. By the proof of the primitive root theorem, $\mathbb{Q}(\alpha_{1}, \alpha_{2}) = \mathbb{Q}(\alpha_{1} + c \alpha_{2})$ for some $c \in \mathbb{Q}^{\times}$ such that $\alpha_{1, i} + c \alpha_{1, j} = \alpha_{1, 1} + c \alpha_{2, 1}$ if and only if $i = j = 1$.

	We claim that $\mathbb{Q}(\sqrt{2}, \sqrt{3}) \cong \mathbb{Q}(\sqrt{2} + \sqrt{3})$ by choosing $c = 1$. Specifically, notice that
	\[
		(\pm \sqrt{2}) + (\pm \sqrt{3}) = \sqrt{2} + \sqrt{3}
	\]
	if and only if we choose the positive sign for both roots on the left hand side. Then $\sqrt{2} +\sqrt{3}$ is a quartic algebraic number with four conjugates given by the four $\pm \sqrt{2} \pm \sqrt{3}$. Now we claim that $\mathbb{Q}(\sqrt{2} + \sqrt{3}, \sqrt{5}) \cong \mathbb{Q}(\sqrt{2} + \sqrt{3} + \sqrt{5})$, again by choosing $c = 1$. w
\end{solution}

\pagebreak

\begin{problem}[Problem 3.]
	A subset $S \subseteq \mathbb{R}$ is called \emph{discrete} if there exists real $\varepsilon > 0$ such that for every two distinct elements $\alpha, \beta \in S$ $\lvert \alpha - \beta \rvert \ge \varepsilon$.

	Let us say that $S$ is \emph{near-discrete} if for every $\alpha \in S$ there exists a real $\varepsilon = \varepsilon(\alpha) > 0$ such that $\lvert \alpha - \beta \rvert \ge \varepsilon$ for every $\beta \in S$ distinct from $\alpha$.
	\begin{enumerate}
		\item Prove that every discrete subset of $\mathbb{R}$ is countable.
		\item Prove that every near-discrete subset of $\mathbb{R}$ is countable.
	\end{enumerate}
\end{problem}

\pagebreak

\begin{problem}[Problem 4.]
	Let $\mu > 0$. We say that $p / q \in \mathbb{Q}$ is a \emph{$\mu$-approximation} to the real number $\alpha$ if $\left \lvert \alpha - p/q \right \rvert \le 1 / q^{\mu}.$ 

	Let $S \subseteq \mathbb{R}$ be a set with the following properties:
	\begin{enumerate}[label = (\arabic*)]
		\item Every element of $S$ has infinitely many rational $3$-approximations
		\item If a rational number $p / q$ is $3$-approximation for some $\alpha \in S$, then it is not a $3$-approximation for any other element of $S$.
	\end{enumerate}

	\begin{enumerate}
		\item Prove that $S$ is countable.
		\item Prove that every $\alpha \in S$ is transcendental.
	\end{enumerate}
\end{problem}

\pagebreak

\begin{problem}[Problem 5.]
	Let $\mathbb{K}$ be a subfield of $\mathbb{C}$. An embedding of $\mathbb{K}$ into $\mathbb{C}$ is an injective field homomorphism $\tau : \mathbb{K} \to \mathbb{C}$.
	\begin{enumerate}
		\item Suppose that $\tau : \mathbb{K} \to \mathbb{C}$ is a field homomorphism such that for some $a \in \mathbb{K}$, $\tau(a) \neq 0$. Prove that $\tau$ is an embedding.

		\item Let $\mathbb{K} = \mathbb{Q}(\alpha)$ for some $\alpha \in \mathbb{C}$. Prove that $\alpha$ is transcendental if and only if there exist infinitely many distinct embeddings of $\mathbb{K}$ into $\mathbb{C}$.
	\end{enumerate}
\end{problem}

\begin{solution}
	\begin{enumerate}
		\item Since $\mathbb{K}$ is a field and  $\ker \tau \subseteq \mathbb{K}$ is an ideal, we must have $\ker \tau = 0$ or $\ker \tau = \mathbb{K}$ (every field has only those two ideals). Since there exists some $a \in \mathbb{K}$ with $\tau(a) \neq 0$ we have $\ker \tau \neq \mathbb{K}$. Thus, $\tau$ has trivial kernel and is injective. That is, it is an embedding $\mathbb{K} \to \mathbb{C}$.

		\item Suppose $\alpha$ is transcendental. We claim any function $\tau : \mathbb{K} \to \mathbb{C}$ fixing $\mathbb{Q}$ and with $\tau(\alpha) \in \mathbb{Q}^{\times} \alpha$ (that is, $\tau(\alpha) = \lambda \alpha$ for some $\lambda \in \mathbb{Q}^\times$), is a ring homomorphism, and thus, an embedding of $\mathbb{K}$. Since $\mathbb{Q}$ is infinite, there are at least infinitely many embeddings $\tau$.

			Suppose $\alpha$ is not transcendental, and instead is algebraic of degree $d$. Then we have shown in class that there are exactly $d$ embeddings $\mathbb{K} \to \mathbb{C}$. (Embeddings of $\mathbb{Q}(\alpha)$ are determined by the image of $\alpha$, embeddings preserve the minimal polynomial of $\alpha$ so $\alpha$ can only be sent to another root of $m_{\alpha}$, of which there are exactly $d$).
	\end{enumerate}
\end{solution}

\pagebreak

\begin{problem}[Problem 6.]
	Let $\alpha, \beta \in \mathbb{C}$, both non-zero, let $\gamma = \alpha / \beta$ and suppose $\mathbb{Q}(\alpha) \cong \mathbb{Q}(\beta)$.
	\begin{enumerate}
		\item Give necessary and sufficient conditions on $\alpha$ and $\beta$ so that $\deg \mathbb{Q}(\gamma) / \mathbb{Q} < \infty$. Prove your answer.

		\item Suppose $\alpha$ and $\beta$ are algebraic and $\gamma \in \mathbb{Q}$. Let
			\[
				f(x) = \sum_{k = 0}^{n} a_{k} x^{k}
			\]
			be the minimal polynomial of $\alpha$. Determine the minimal polynomial $g(x)$ of $\beta$.

		\item Suppose $\alpha$ and $\beta$ are algebraic and $f(x) = g(x)$. Is it true that $\gamma \in \mathbb{Q}$? Prove or give a counterexample.
	\end{enumerate}
\end{problem}

\pagebreak

\begin{problem}[Problem 7.]
	Let $\{ n_{j} \}_{j = 1}^{\infty}$ be a sequence of natural numbers such that
	\[
		\lim_{j \to \infty} \frac{n_{j + 1}}{n_{j}} = \infty.
	\]
	Define
	\[
		f(x) = \sum_{j = 1}^{\infty} x^{n_{j}},
	\]
	and let $\alpha$ be a rational number such that $0 < \alpha < 1$. Prove that $f(\alpha)$ is transcendental.
\end{problem}

\pagebreak

\begin{problem}[Problem 8.]
	Let $\alpha_{1}, \dots, \alpha_{n}$ be algebraic integers. Give necessary and sufficient conditions on these numbers for the polynomial
	\[
		f(x) = (x - \alpha_{1}) \cdots (x - \alpha_{n})
	\]
	to have rational integer coefficients. Prove your answer.
\end{problem}

\begin{solution}
	$f$ has rational integer coefficients if and only if the set $\{ \alpha_{1}, \dots, \alpha_{n} \}$ can be partitioned such that the complete set 
\end{solution}

\end{document}
